% Einheit 4 von 4 – Sachaufgaben (bank/quadratische-gleichungen/e4.jsonl, zusammenbau v0.1)
%% TODO Zweigzeile Teil 1: was hier gelernt wird (Satz in Schülersprache aus dem Einheitstitel) – Einheit 4
\einheitenkopf[e4]{Einheit 4 von 4 · Sachaufgaben}
\zweigzeile{ab Kl. 9, je nach Buch bis Kl. 10 · keine P10-Aufgabe}
\setcounter{aufgabe}{20}

%% TODO Titel: Ich-kann-Satz fehlt in der Bank (Platzhalter: Kettenname) – Nr. 21
%% TODO Anweisung: ein Satz über der ersten Teilaufgabe fehlt in der Bank – Nr. 21
\begin{aufgabe}{Sachaufgaben}
%% TODO \rechenplatz{n}: Schrittzahl des Grundfalls fehlt in der Bank (nur bei fester Schreibform, 2.3 b)
\begin{teile}
\teil Zahlenrätsel „Das Quadrat einer Zahl ist $25$“: Lösungen $5$ und $-5$ – welche passen?\\ \kreuz{nur $5$} \\ \kreuz{nur $-5$} \\ \kreuz{beide}
\teil Das Quadrat einer Zahl $x$ ist $144$: $x^2 = 144$. Welche Zahlen? x1 = \leerfeld , x2 = \leerfeld
\teil Vermehrt man das Quadrat einer Zahl $x$ um $7$, erhält man $88$. Welche Zahlen? Gleichung: \leerfeld ; x1 = \leerfeld , x2 = \leerfeld
\teil Das Produkt einer Zahl $x$ mit ihrem Nachfolger ist $56$. Welche Zahlen? x1 = \leerfeld , x2 = \leerfeld
\teil Ein rechteckiger Garten ist $5$ m länger als breit und hat die Fläche $84\,\text{m}^2$. Skizze, Gleichung mit der Breite $x$, Lösungen? Gleichung: \leerfeld ; x1 = \leerfeld , x2 = \leerfeld
\teil Ein Rechteck ist $4$ cm länger als breit und hat die Fläche $96\,\text{cm}^2$. Wie breit ist es? Begründe, welche Lösung wegfällt. Breite: \leerfeld[cm]
\teil Ein rechteckiger Tisch ist $40$ cm länger als breit und hat die Fläche $6000\,\text{cm}^2$. Breite und Länge? Antworte im Satz. Antwort: \leerfeld
\teil Lukas sagt: Ein Rechteck, das $3$ cm länger als breit ist und $70\,\text{cm}^2$ Fläche hat, ist $7$ cm breit. Stimmt das? \janein
\teil Ein Rechteck ist $4$ m länger als breit. Verlängert man beide Seiten um $2$ m, beträgt die Fläche $96\,\text{m}^2$. Wie breit und wie lang ist das ursprüngliche Rechteck? Breite: \leerfeld[m] , Länge: \leerfeld[m]
\end{teile}
\end{aufgabe}

%% TODO Titel: Ich-kann-Satz fehlt in der Bank (Platzhalter: Kettenname) – Nr. 22
%% TODO Anweisung: ein Satz über der ersten Teilaufgabe fehlt in der Bank – Nr. 22
\begin{aufgabe}{Quadrat mit Rand}
\begin{teile}
\teil Ein quadratisches Bild bekommt ringsum einen $5$ cm breiten Rahmen. Mit Rahmen ist die Fläche $900\,\text{cm}^2$. Wie lang ist die Seite des Bildes? \leerfeld[cm]
\end{teile}
\end{aufgabe}

%% TODO Titel: Ich-kann-Satz fehlt in der Bank (Platzhalter: Kettenname) – Nr. 23
%% TODO Anweisung: ein Satz über der ersten Teilaufgabe fehlt in der Bank – Nr. 23
\begin{aufgabe}{Sachaufgaben – Fehler finden, Begründen}
\begin{teile}
\teil Ein Rechteck ist $2$ m länger als breit und hat $35\,\text{m}^2$ Fläche. Wo ist der Fehler? \rechnung{x \cdot (x + 2) &= 35 \\ x^2 + 2x - 35 &= 0 \\ x_1 = 5, \quad x_2 &= -7} Antwort: Die Breite ist $-7$ m.
\teil Für die Breite eines Rechtecks ergibt die Gleichung $x_1 = 9$ und $x_2 = -13$. Begründe, warum nur eine Lösung die Breite ist.
\end{teile}
\end{aufgabe}
